\chapter*{Preface}
The fourteen tomes before this one prove theorems. This one is about what happens after the last theorem of a book: how a
mathematician reads a research paper, tests an idea on a computer, trusts a heuristic, checks a proof, writes it down, and
chooses the next problem. It is a workshop, and every tool in it is used on real mathematics, proved in full.

The first part looks at research from the inside. Chapter~1 follows one problem through its whole life: Erd\H os's question,
from 1931, on how small a set of integers with distinct subset sums can be, with its bounds, its constructions, its prizes,
and its disproof in 2026, whose pieces we check ourselves. Chapter~2 teaches the three readings of a paper on the two-page
note of Dubroff, Fox and Xu about the same problem. Along the way it finds a tempting simplification that fails, which is a
lesson in itself.

The second part is four guided readings of papers that made probability and number theory meet. Bayer and Diaconis's
``seven shuffles'' theorem is proved here in its sharp form, with the limiting shape of the cutoff. Radziwiłł and
Soundararajan's short proof of Selberg's central limit theorem for the zeta function is mapped onto the series and extended
to several points at once, where a log-correlated structure appears. Wood's theory of random groups is read through its simplest
case, the cokernels of random symmetric matrices, and tested on the sandpile groups of random graphs. Harper's theorem that a random multiplicative function has ``better than square
root cancellation'' is explained through a branching random walk, in which the heart of the matter, a ballot theorem below
a barrier, is proved completely.

The third part is the workshop proper. Chapter~7 is experimental mathematics: digits of $\pi$ computed from the middle,
and integer relations discovered by lattice reduction, with the LLL algorithm proved. Chapter~8 is the art of the heuristic,
on sums of three cubes, from the congruence obstruction modulo $9$ to the local densities and to the factorisation that
made the search for $33$ and $42$ possible. Chapter~9 explains proof assistants, shows Lean at work, and proves the
Curry--Howard correspondence that makes them possible. Chapter~10 is about writing mathematics, with Jordan's theorem for
polygons as a worked example of a proof that is easy to believe and must be written with care. Chapter~11 offers seven open
problems, each with a partial result proved in full, and Chapter~12 five projects that start with a computation and end
with a question nobody has answered.

The book assumes the series: it never reproves a result of an earlier tome, but says precisely where each one is proved.
Results admitted for culture are labelled as remarks, with references. Each chapter opens with the story behind its
ideas, checked against the sources, and ends with exercises in the style of the series: a hint, then a solution that first
explains how to find the idea. Every chapter has a Python notebook that recomputes its numbers and opens a window on
computer science: lattice attacks on knapsack cryptosystems, point-in-polygon algorithms, program synthesis by proof
search, segmented sieves and backtracking.

Research is the place where the exercises run out. The pleasure of mathematics is not diminished there; it is at its
purest, because for the first time nobody knows the answer. This book tries to show what that feels like, and to give
the reader everything needed to start.
